% TOP_PROBLEM: {"problemId": "problem.iitaka-conjecture-c-n-m-on-subadditivity-of-kodaira-dimension", "canonicalTitle": "Iitaka conjecture C_{n,m} on subadditivity of Kodaira dimension", "releaseRank": 114, "primaryDomain": "geometry_topology", "primaryDomainLabel": "Geometry & topology", "displayStatus": "open_with_solved_subcases", "releaseStatus": "open_with_solved_subcases"}
% !TeX program = lualatex
% Definition and sources only; this file is not an LLM solution attempt.
\documentclass[11pt]{article}
\usepackage[margin=25mm]{geometry}
\usepackage{fontspec}
\setmainfont{Latin Modern Roman}
\usepackage{amsmath,amssymb,mathtools}
\usepackage{unicode-math}
\setmathfont{Latin Modern Math}
\usepackage{xurl,hyperref}
\hypersetup{colorlinks=true,urlcolor=blue}
\setcounter{secnumdepth}{0}
\setlength{\parindent}{0pt}
\setlength{\parskip}{5pt}
\setlength{\emergencystretch}{3em}
\begin{document}
\section*{\texorpdfstring{Iitaka conjecture C\_\{n,m\} on subadditivity of Kodaira dimension}{Iitaka conjecture C\_\{n,m\} on subadditivity of Kodaira dimension}}\label{114}
\subsection{Definitions and mathematical statement}
For a smooth projective variety $Z$, its Kodaira dimension $\kappa(Z)$ is the growth exponent of the spaces $H^0(Z,mK_Z)$, with value $-\infty$ if all positive pluricanonical spaces vanish. For the surjective morphism and general fiber in the record, the target is
\[
 \kappa(X)\ge\kappa(Y)+\kappa(F).
\]
The usual fibration convention uses connected fibers and the characteristic-zero setting of the cited source; the catalogue's abbreviated sentence does not separately spell out those conventions. When one term on the right is $-\infty$, the inequality is interpreted in the extended ordered sense.
\par\smallskip\textit{Formulation and concept references:} \href{https://arxiv.org/abs/0806.4413}{[S1]}.

\subsection{Short English statement}
Is the Kodaira dimension of a projective fibration at least the sum of those of its base and general fiber?
\subsection{Sources}
\begin{itemize}
\item[S1]C. Birkar, Iitaka conjecture C\_\{n,m\} in dimension six, arXiv:0806.4413 [math.AG], 2008.
\url{https://arxiv.org/abs/0806.4413}.
\textit{Formulation source.}
\end{itemize}
\end{document}
